\documentclass[10pt,english]{article}

\usepackage{amsfonts,url,tikz,amsmath}

% Paper setup
\evensidemargin=0in
\oddsidemargin=0in
\textwidth=6.25in
\topmargin=-0.5in
\headheight=0.0in
\headsep=0.5in
\textheight=9.0in
\footskip=0.5in

\input{institution}

\newcommand{\Interior}[1]{\ensuremath{{#1}^{\circ}}}
\newcommand{\Closure}[1]{\ensuremath{\overline{#1}}}
\newcommand{\Complement}[1]{\ensuremath{{#1}^{c}}}

\newcommand{\Expect}{\ensuremath{\mathrm{E}}}
\newcommand{\vecnot}{\underline}
\newcommand{\RealNumbers}{\ensuremath{\mathbb{R}}}
\newcommand{\RationalNumbers}{\mathbb{Q}}
\newcommand{\ComplexNumbers}{\mathbb{C}}
\newcommand{\Real}{\mathrm{Re}}
\newcommand{\Span}{\mathrm{span}}
\newcommand{\Rank}{\mathrm{rank}}
\newcommand{\Nullity}{\mathrm{nullity}}
\newcommand{\Trace}{\mathrm{tr}}
\newcommand{\Diag}{\mathrm{diag}}
\newcommand{\dd}{\mathrm{d}}
\DeclareMathOperator*{\esssup}{ess\,sup}

% Use < , > inner product
\newcommand{\inner}[2]{{\left\langle #1 \mskip2mu , #2 \right\rangle}}
\newcommand{\tinner}[2]{{\langle #1 \mskip1mu , #2 \rangle}}

% Use < | > inner product
%\newcommand{\inner}[2]{{\left\langle #1 \mskip2mu \middle| \mskip2mu #2 \right\rangle}}
%\newcommand{\tinner}[2]{{\langle #1 \mskip1mu | \mskip1mu  #2 \rangle}}



\begin{document}

\title{ECE 586: Vector Space Methods \\ Lecture 16: Frequently Asked Questions}
\author{Henry D. Pfister \\ \univ}
\date{October 18, 2026}

\maketitle

\section{Derivatives as Linear Maps}

\paragraph{Why is a derivative a linear transformation rather than just a number?}
For $f\colon\mathbb{R}^n\to\mathbb{R}^m$, the derivative at $x$ maps a perturbation $h$ to the first-order output change $J(x)h$.
The Jacobian is an $m\times n$ matrix representation of that linear map.
For $n=m=1$, a linear map is multiplication by a scalar, which recovers the familiar derivative (slide 2, ``Derivatives'').

\paragraph{What does the Fr\'echet derivative definition require?}
It requires a bounded linear map $T$ for which
\[
f(x+h)=f(x)+Th+r(h),\qquad \frac{\|r(h)\|}{\|h\|}\longrightarrow0.
\]
The same $T$ must work as $h$ tends to zero in every direction, not just along a fixed line.
This is the definition on slide 5, ``What is Meant by Differentiable?''

\paragraph{What does a bound on the derivative tell us about the function?}
If $f$ is differentiable on a neighborhood of a convex set $A$ and $\|f'(x)\|_{\mathrm{op}}\leq L$ throughout $A$, then $\|f(y)-f(x)\|\leq L\|y-x\|$ for $x,y\in A$.
Convexity lets us follow the entire line segment between the points while using the derivative bound.
The chain rule on slide 6 composes derivative maps in the same order as the corresponding functions.

\paragraph{How do we compute a directional derivative in the slide example?}
Substitute $x+th$ into the function and differentiate with respect to the scalar $t$ at zero.
For $f(x_1,x_2)=(x_1x_2,x_1+x_2^2)$, $x=(1,1)$, and $h=(1,2)$, this gives $(3,5)$, as on slide 7.
When $f$ is Fr\'echet differentiable, the answer equals $f'(x)h$.

\paragraph{Does existence of every directional derivative imply Fr\'echet differentiability?}
No. Directional derivatives test one direction at a time, and can fail even to depend linearly on the direction.
The slides call the function G\^ateaux differentiable when these derivatives define a continuous linear function of the direction.
Even that condition is weaker than the uniform first-order remainder required for Fr\'echet differentiability.

\section{Gradients and Descent}

\paragraph{Why does a negative directional derivative give a descent direction?}
If $\delta f(x;h)=a<0$, then $[f(x+th)-f(x)]/t$ is negative for all sufficiently small positive $t$.
Multiplying by $t>0$ gives $f(x+th)<f(x)$.
This is a local statement; it does not say that every positive step size will decrease the function (slide 8).

\paragraph{What is the distinction between the derivative and the gradient?}
For a real-valued function on a real Hilbert space, $f'(x)$ is a bounded linear functional, whereas $\nabla f(x)$ is a vector in the space.
The Riesz theorem identifies them through $f'(x)h=\inner{h}{\nabla f(x)}$.
In Euclidean coordinates, the derivative is a row matrix and the gradient is its transpose (slide 9).
A general Banach space has no canonical inner-product identification of its dual with itself.

\paragraph{Why does gradient descent use the negative gradient?}
Among unit directions in a Hilbert space, $-\nabla f(x)/\|\nabla f(x)\|$ gives the most negative directional derivative when the gradient is nonzero.
This follows from Cauchy--Schwarz.
The iteration on slide 10 is $x_{n+1}=x_n-\delta_n\nabla f(x_n)$; the step size still needs to be chosen appropriately.

\paragraph{How does a Lipschitz gradient give a useful step-size bound?}
Slide 11 bounds the first-order approximation error by $L\|y-x\|^2/2$.
Substituting $y=x-\delta\nabla f(x)$ yields
\[
f(y)\leq f(x)-\delta\left(1-\frac{L\delta}{2}\right)\|\nabla f(x)\|^2.
\]
For $L>0$, a step $0<\delta<2/L$ therefore decreases $f$ when the gradient is nonzero, provided the bound applies along that step.

\paragraph{Does a zero gradient guarantee a global minimum?}
No. For a general differentiable function, it may indicate a local maximum or saddle point.
For a differentiable convex function on the whole space, a zero gradient is sufficient for a global minimum.
Convexity and constrained optimality are developed in Videos 21 and 23.

\end{document}
